Bagian ini mensintesis temuan utama dari studi yang ditinjau dengan tujuan 
mengidentifikasi pola asumsi perilaku lawan serta implikasinya terhadap 
desain metodologis pemodelan lawan dalam \textit{repeated games}. 
Sintesis difokuskan pada dimensi perilaku yang dapat diamati dari riwayat 
interaksi, alih-alih pada asumsi internal seperti struktur optimisasi reward 
atau arsitektur pembelajaran yang digunakan.

Pendekatan ini penting karena dalam banyak karya terdahulu, pemodelan 
perilaku lawan sering kali tidak dipisahkan secara konseptual dari proses 
optimisasi kebijakan agen. Akibatnya, representasi perilaku lawan terenkapsulasi 
di dalam parameter kebijakan atau fungsi nilai, sehingga sulit untuk 
mengevaluasi kualitas prediksi perilaku secara independen dari reward yang 
dicapai. Untuk mengidentifikasi celah tersebut secara sistematis, tinjauan ini 
mengabstraksikan asumsi perilaku lawan pada tingkat eksternal yang dapat 
diinferensi dari \textit{interaction traces}.

Tabel~\ref{tab:ringkasan_terkait} merangkum berbagai pendekatan dalam \textit{opponent modelling} dan pembelajaran multi-agen pada interaksi strategis berulang. Pendekatan berbasis pemodelan eksplisit terhadap lawan ditunjukkan oleh Online Opponent Modeling (O2M) yang dilaporkan mampu beradaptasi lebih baik dan memperoleh reward rata-rata lebih tinggi dibandingkan baseline (\cite{qiao_o2m_2024}). Dalam konteks dinamika evolusioner, strategi Zero Determinant menunjukkan proporsi kerja sama akhir yang lebih tinggi dibandingkan strategi Extort, mengindikasikan peran signifikan dinamika populasi dalam membentuk hasil kolektif (\cite{zhu_evolutionary_2025}).

Pendekatan berbasis pembelajaran mendalam juga menunjukkan kemampuan adaptif terhadap dinamika lawan. Hierarchical Gifting DQN mampu mendeteksi perubahan strategi secara real-time dan menyesuaikan insentif kerja sama secara dinamis (\cite{lv_inducing_2023}). Model berbasis LSTM dilaporkan dapat mengaproksimasi dinamika no-regret secara akurat dan mengeksploitasi lawan secara menguntungkan pada tingkat non-stasioneritas rendah (\cite{li_exploiting_2025}). Selain itu, pendekatan Deep Multiagent Reinforcement Learning pada Social Prisoner's Dilemma mampu mencapai kerja sama mutual dan beradaptasi terhadap perubahan strategi lawan (\cite{wang_achieving_2019}).

Beberapa studi menekankan pentingnya struktur interaksi dan kopling dinamika. Environment-Strategy Coupling Model menunjukkan bahwa kopling antara umpan balik lingkungan dan dinamika strategi meningkatkan stabilitas kerja sama kelompok (\cite{di_coupling_2023}), sementara Extended Social Particle Swarm Model memperlihatkan bahwa ukuran populasi dan laju pembaruan informasi memengaruhi pola kerja sama yang muncul (\cite{elhamer_effects_2020}). Pendekatan Teamwork Game dengan MA-MAB bahkan mampu mereproduksi pola perilaku mirip manusia seperti \textit{social loafing} dan \textit{compensation} (\cite{gomez_grounded_2025}).

Dalam konteks pemodelan internal dan kesadaran pembelajaran, Internal Model dilaporkan mencapai konvergensi lebih cepat dibanding TD-learning standar, meskipun performanya menurun ketika perilaku lawan bergantung langsung pada aksi agen (\cite{freire_modeling_2023}). Symmetric Learning Awareness (SLA) meningkatkan stabilitas keseimbangan dan mengurangi siklikitas pada metode gradien standar melalui pemodelan eksplisit terhadap proses pembelajaran lawan (\cite{hu_modeling_2023}). Pendekatan Higher-order Theory of Mind (ToM $k$) menunjukkan bahwa peningkatan level ToM memberikan pengaruh signifikan hingga tingkat tertentu, dengan \textit{diminishing returns} setelah ToM-3 (\cite{de_weerd_higher-order_2022}). Sementara itu, mekanisme Suggestion Sharing meningkatkan kerja sama dibandingkan skema berbagi nilai atau kebijakan (\cite{jin_achieving_2025}), dan pendekatan Ensemble-Training Cooperative Agent menunjukkan peningkatan robustnes dan generalisasi meskipun dengan biaya komputasi yang lebih tinggi (\cite{perera_learning_2025}).

Secara keseluruhan, literatur menunjukkan keberagaman pendekatan dalam memodelkan dan merespons dinamika perilaku lawan, baik melalui pembelajaran eksplisit, mekanisme kopling lingkungan, maupun adaptasi berbasis populasi. 

\section{Asumsi Perilaku Lawan dalam pemodelan lawan untuk \textit{Repeated Games}}

Untuk mengoperasionalkan asumsi perilaku lawan secara konsisten pada 
pengaturan permainan yang beragam, tinjauan ini mengkarakterisasi perilaku 
berdasarkan dependensi yang dapat diamati dari riwayat interaksi. 
Secara khusus, aksi lawan dianalisis berdasarkan apakah ia:

\begin{enumerate}
    \item Bervariasi lintas kondisi lingkungan dalam permainan yang sama,
    \item Dimediasi oleh dinamika tingkat populasi,
    \item Merespons secara langsung aksi agen,
    \item Menunjukkan divergensi aksi pada riwayat interaksi terkini yang ekuivalen.
\end{enumerate}

Kriteria-kriteria ini sepenuhnya berbasis observasi eksternal dan tidak 
mengasumsikan akses terhadap model internal, mekanisme pembelajaran, 
maupun fungsi objektif lawan. Dengan demikian, kategorisasi ini bukan 
taksonomi formal algoritma, melainkan kerangka analitis untuk memungkinkan 
perbandingan lintas paradigma pembelajaran.

Dalam studi yang tidak menyatakan asumsi perilaku secara eksplisit, 
klasifikasi diturunkan secara konservatif dari desain eksperimen dan dinamika 
interaksi yang dilaporkan. Hasil kategorisasi dirangkum pada 
Tabel~\ref{tab:opponent-behavior}.

Distribusi asumsi dependensi perilaku lawan dalam literatur menunjukkan variasi tingkat kompleksitas interaksi yang dimodelkan. Studi pada \cite{jin_achieving_2025} berfokus pada respons langsung terhadap aksi agen tanpa mempertimbangkan efek populasi maupun variasi lintas lingkungan. Pendekatan pada \cite{gomez_grounded_2025} menempatkan perilaku dalam konteks populasi, namun tanpa asumsi respons langsung terhadap agen atau divergensi pada riwayat yang ekuivalen. Sementara itu, \cite{elhamer_effects_2020} menggabungkan dinamika populasi dan respons agen, tetapi tidak secara eksplisit memasukkan asumsi divergensi perilaku pada riwayat identik.

Sebaliknya, sejumlah studi secara eksplisit mengasumsikan adanya respons agen sekaligus potensi divergensi perilaku. Pendekatan pada \cite{qiao_o2m_2024}, \cite{lv_inducing_2023}, \cite{li_exploiting_2025}, \cite{freire_modeling_2023}, \cite{hu_modeling_2023}, \cite{wang_achieving_2019}, dan \cite{de_weerd_higher-order_2022} memodelkan lawan yang secara langsung merespons aksi agen serta memungkinkan terjadinya divergensi aksi meskipun riwayat interaksi terkini setara. Studi pada \cite{perera_learning_2025} memperluas asumsi ini dengan memasukkan dimensi populasi, sehingga perilaku lawan dipengaruhi baik oleh interaksi agen maupun struktur populasi, serta tetap memungkinkan divergensi. Pada tingkat kompleksitas yang lebih tinggi, \cite{zhu_evolutionary_2025} dan \cite{di_coupling_2023} mempertimbangkan variasi lintas lingkungan sekaligus dinamika populasi, dengan potensi divergensi perilaku yang muncul dari interaksi struktural dan evolusioner.


\begin{table}[!htbp]
\centering
\caption{Asumsi perilaku lawan berdasarkan dependensi perilaku yang dapat diamati.}\label{tab:opponent-behavior}
\renewcommand{\arraystretch}{1.1}
\setlength{\tabcolsep}{4pt}

\begin{tabularx}{\textwidth}{>{\raggedright\arraybackslash}X cccc}
\toprule
\textbf{Referensi}
& \textbf{Lingkungan} 
& \textbf{Populasi} 
& \textbf{Agen}
& \textbf{divergensi} \\
\midrule

~\cite{jin_achieving_2025} 
& --- & --- & \checkmark{} & --- \\

\addlinespace{}

~\cite{gomez_grounded_2025} 
& --- & \checkmark{} & --- & --- \\

\addlinespace{}

~\cite{elhamer_effects_2020} 
& --- & \checkmark{} & \checkmark{} & --- \\

\addlinespace{}

~\cite{qiao_o2m_2024}
& --- & --- & \checkmark{} & \checkmark{} \\

\addlinespace{}

~\cite{lv_inducing_2023}
& --- & --- & \checkmark{} & \checkmark{} \\

\addlinespace{}

~\cite{li_exploiting_2025}
& --- & --- & \checkmark{} & \checkmark{} \\

\addlinespace{}

~\cite{freire_modeling_2023}
& --- & --- & \checkmark{} & \checkmark{} \\

\addlinespace{}

~\cite{hu_modeling_2023}
& --- & --- & \checkmark{} & \checkmark{} \\

\addlinespace{}

~\cite{wang_achieving_2019}
& --- & --- & \checkmark{} & \checkmark{} \\

\addlinespace{}

~\cite{de_weerd_higher-order_2022}
& --- & --- & \checkmark{} & \checkmark{} \\

\addlinespace{}

~\cite{perera_learning_2025}
& --- & \checkmark{} & \checkmark{} & \checkmark{} \\

\addlinespace{}

~\cite{zhu_evolutionary_2025}
& \checkmark{} & \checkmark{} & --- & \checkmark{} \\

\addlinespace{}

~\cite{di_coupling_2023}
& \checkmark{} & \checkmark{} & \checkmark{} & \checkmark{} \\

\bottomrule
\end{tabularx}

\begin{flushleft}
\vspace{2pt}
\footnotesize
\noindent~\textit{Catatan:} \\
Lingkungan — perilaku bervariasi lintas lingkungan dalam permainan yang sama;\\
Populasi — perilaku dimediasi oleh interaksi tingkat populasi;\\
Agen — perilaku merespons secara langsung aksi agen;\\
Divergen — divergensi aksi terjadi pada riwayat interaksi terkini yang ekuivalen.\\
Tanda centang menunjukkan adanya dependensi.
\end{flushleft}
\end{table}

Sejumlah pola konsisten muncul dari Tabel~\ref{tab:opponent-behavior}. 
Mayoritas studi terkini mengasumsikan lawan yang responsif terhadap aksi agen 
serta menunjukkan divergensi perilaku pada riwayat interaksi yang setara. 
Hal ini mencerminkan pergeseran fokus riset dari optimisasi terhadap strategi 
lawan yang tetap menuju interaksi dengan lawan yang adaptif dan dinamis.

\section{Pendekatan metodologis dalam pemodelan lawan untuk \textit{Repeated Games}}

Distribusi asumsi perilaku lawan pada Tabel~\ref{tab:opponent-behavior} mencerminkan pergeseran metodologis yang lebih luas dalam riset pemodelan lawan, dari pengaturan dengan asumsi lawan yang tetap dan terkontrol menuju lawan yang adaptif, heterogen, dan berperilaku tidak stasioner. Namun, karakterisasi berbasis asumsi perilaku semata belum menjelaskan bagaimana kompleksitas tersebut dihadapi secara komputasional di sisi agen. Untuk itu, Tabel~\ref{tab:klasifikasi_pemodelan_lawan} mereorganisasi studi-studi terdahulu berdasarkan paradigma pemodelan dominan dan mekanisme pembelajaran yang digunakan oleh agen.

Pendekatan-pendekatan tersebut dapat dibedakan lebih lanjut berdasarkan apakah adaptasi terhadap lawan ditangani secara eksplisit atau implisit. Paradigma pemodelan lawan yang eksplisit, seperti \textit{gradient-based opponent shaping}~\cite{qiao_o2m_2024, hu_modeling_2023, wang_achieving_2019}, \textit{recursive belief reasoning}~\cite{freire_modeling_2023, de_weerd_higher-order_2022}, serta \textit{system identification}~\cite{li_exploiting_2025} secara langsung membangun representasi internal perilaku lawan untuk memprediksi atau memengaruhi respons lawan di masa depan. Sebaliknya, pendekatan implisit, termasuk \textit{reactive reinforcement learning}~\cite{lv_inducing_2023}, \textit{population-based training}~\cite{perera_learning_2025}, dinamika evolusioner~\cite{zhu_evolutionary_2025, di_coupling_2023, elhamer_effects_2020}, dan \textit{bandit-based learning-in-games}~\cite{gomez_grounded_2025}, beradaptasi terhadap lawan tanpa mempertahankan model perilaku lawan yang terpisah.

Meskipun demikian, baik pendekatan eksplisit maupun implisit pada umumnya merumuskan dan mengevaluasi adaptasi terhadap lawan hanya pada tingkat respons satu langkah (\textit{one-step interaction}), tanpa melakukan pemodelan atau evaluasi eksplisit terhadap konsekuensi perilaku lawan dalam horizon multi-langkah.

\begin{table}[ht]
\centering
\caption{Klasifikasi Pendekatan Pemodelan Lawan}
\label{tab:klasifikasi_pemodelan_lawan}
\begin{tabular}{p{2cm} p{2.0cm} p{4cm} p{4cm}}
\toprule
\textbf{Ref} & \textbf{Kategori} & \textbf{Pendekatan} & \textbf{Model} \\
\midrule

\cite{wang_achieving_2019} 
& Implisit 
& Bandit-based Learning 
& Multi-Armed Bandit \\

\cite{gomez_grounded_2025} 
& Implisit 
& Bandit-based Learning 
& Multi-Armed Bandit (MA-MAB) \\

\cite{zhu_evolutionary_2025} 
& Implisit 
& Dinamika Evolusioner 
& Strategi berbasis fitness \\

\cite{elhamer_effects_2020} 
& Implisit 
& Dinamika Evolusioner 
& Extended Social Particle Swarm (SPS) \\

\cite{perera_learning_2025} 
& Implisit 
& Population-based Training 
& MARL berbasis ensemble \\

\cite{wang_achieving_2019} 
& Implisit 
& Reactive Reinforcement Learning 
& Deep RL (DQN) \\

\cite{lv_inducing_2023} 
& Eksplisit 
& Gradient-based Opponent Shaping 
& Pembaruan sadar lawan \\

\cite{hu_modeling_2023} 
& Eksplisit 
& Gradient-based Opponent Shaping 
& Symmetric Learning Awareness (SLA) \\

\cite{de_weerd_higher-order_2022} 
& Eksplisit 
& Recursive Belief Reasoning 
& Higher-order Theory of Mind \\

\cite{li_exploiting_2025} 
& Eksplisit 
& System Identification 
& NARX / LSTM Strategy \\

\bottomrule
\end{tabular}
\end{table}

Kebanyakan Pendekatan yang memperlakukan perilaku lawan sebagai proses sekuensial yang dapat 
dipelajari, serta dievaluasi berdasarkan 
kualitas prediksi terhadap lintasan aksi di masa depan, bukan semata-mata 
berdasarkan payoff agregat jangka panjang. Pendekatan semacam ini secara 
alami menekankan pentingnya \textit{forecasting} multi-langkah dalam 
menghadapi dinamika interaksi berulang, terutama ketika keputusan pada 
fase awal interaksi memiliki implikasi strategis terhadap lintasan berikutnya.

Dalam konteks ini, pendekatan berbasis model sekuensial memori seperti \textit{Long Short-Term Memory} (LSTM) menawarkan alternatif metodologis yang berbeda. LSTM dapat dipandang sebagai pendekatan yang memodelkan perilaku lawan sebagai proses dinamis berbasis riwayat interaksi. 

Model LSTM dilatih untuk meminimalkan kesalahan prediksi perilaku lawan tanpa melibatkan sinyal reward agen. Dengan demikian, kualitas model lawan dapat dievaluasi secara independen menggunakan metrik prediksi sekuensial, sebelum diintegrasikan ke dalam proses pengambilan keputusan.

Selain itu, struktur rekuren pada LSTM memungkinkan forecast multi-langkah (\textit{multi-step forecasting}) terhadap aksi lawan memanfaatkan skema prediksi rekursif atau \textit{rolling horizon}. Kemampuan ini relevan dalam \textit{repeated games}, di mana keputusan optimal agen sering kali bergantung pada ekspektasi terhadap respons lawan tidak hanya pada langkah berikutnya, tetapi juga pada beberapa langkah interaksi selanjutnya. Dengan menyediakan estimasi lintasan perilaku lawan pada horizon terbatas, agen dapat melakukan evaluasi strategi secara lebih prospektif dibandingkan pendekatan satu-langkah.

\section{Bagaimana efektivitas strategi pemodelan lawan dievaluasi dalam \textit{Repeated Games}?}

Praktik evaluasi secara implisit mendefinisikan apa yang dianggap sebagai keberhasilan dalam interaksi multi-agent yang adaptif, baik dalam bentuk hasil kerja sama, ketahanan terhadap eksploitasi, stabilitas perilaku, maupun akurasi prediksi. Tabel~\ref{tab:evaluation_metrics} merangkum lingkungan evaluasi dan metrik yang digunakan dalam penelitian-penelitian terdahulu, dengan tujuan mengidentifikasi pola evaluasi yang berulang serta aspek-aspek yang relatif terabaikan.

Berbagai penelitian terdahulu mengevaluasi pendekatan pemodelan lawan dalam lingkungan dan metrik yang beragam.~(\cite{qiao_o2m_2024}) mengevaluasi metode mereka dalam pengaturan \textit{self-play} simetris dengan fokus pada eror prediksi (MSE) selama pelatihan \textit{offline} serta akurasi representasi memori laten.~(\cite{zhu_evolutionary_2025}) menggunakan jaringan \textit{scale-free} dengan strategi \textit{zero-determinant} dan menilai dinamika melalui frekuensi kerja sama dan eksploitasi. Pendekatan \textit{reactive reinforcement learning} yang diajukan ~(\cite{lv_inducing_2023}) dievaluasi terhadap lawan adaptif dengan metrik utama berupa nilai reward yang dirata-ratakan pada beberapa tipe lawan.~(\cite{gomez_grounded_2025}) menguji pendekatan mereka dalam simulator \textit{teamwork-game} khusus, menggunakan produktivitas tim agregat, konvergensi menuju \textit{Nash equilibrium}, serta kontribusi individu sebagai indikator kinerja.

Dalam konteks dinamika evolusioner~(\cite{di_coupling_2023}) dan~(\cite{elhamer_effects_2020}) memanfaatkan simulator permainan evolusioner dan ruang kontinu, dengan metrik seperti tingkat kerja sama, fraksi kooperator, serta stabilitas klaster.~(\cite{perera_learning_2025}) mengevaluasi pendekatan berbasis populasi dalam \textit{repeated matrix games} menggunakan populasi sintetis, dengan fokus pada payoff rata-rata, robustness, dan generalisasi. ~(\cite{li_exploiting_2025}) pada \textit{repeated zero-sum games} mengukur galat prediksi dan payoff kumulatif sebagai indikator efektivitas eksploitasi lawan. (\cite{freire_modeling_2023}) dan ~(\cite{hu_modeling_2023}) mengkaji \textit{repeated matrix games} dengan menilai efektivitas, stabilitas, akurasi prediksi, payoff rata-rata, serta kecepatan konvergensi.

Selain itu ~(\cite{wang_achieving_2019}) mengembangkan lingkungan SPD dua dimensi khusus dan mengevaluasi reward rata-rata serta kesejahteraan sosial, sementara ~(\cite{jin_achieving_2025}) memanfaatkan benchmark MARL dengan metrik return ternormalisasi dan konvergensi.~(\cite{de_weerd_higher-order_2022}) menggunakan simulasi \textit{Colored Trails} dan menilai performa melalui skor agen serta kesejahteraan sosial. Secara keseluruhan, variasi lingkungan dan metrik ini menunjukkan bahwa evaluasi dalam literatur cenderung berfokus pada indikator payoff, konvergensi, atau tingkat kerja sama, sementara pengukuran eksplisit terhadap kualitas peramalan trajektori perilaku lawan dalam horizon multi-langkah masih relatif terbatas.

\begin{table}[!htbp]
\centering
\caption{Lingkungan evaluasi dan metrik dalam penelitian}\label{tab:evaluation_metrics}
\renewcommand{\arraystretch}{1.2}
\setlength{\tabcolsep}{4pt}
\begin{tabularx}{\columnwidth}{c >{\raggedright\arraybackslash}X >{\raggedright\arraybackslash}X}
\midrule
\textbf{Ref.} & \textbf{Lingkungan Evaluasi} & \textbf{Metrik} \\
\midrule
~\cite{qiao_o2m_2024} & Self-play simetris & MSE selama pelatihan offline; akurasi memori laten \\
~\cite{zhu_evolutionary_2025} & Jaringan scale-free dengan strategi zero-determinant & Frekuensi kerja sama (C) dan eksploitasi (E) \\
~\cite{lv_inducing_2023} & Opponent adaptif (dirata-ratakan pada beberapa opponent) & Nilai reward \\
~\cite{gomez_grounded_2025} & Simulator teamwork-game khusus & Produktivitas tim agregat; konvergensi ke Nash equilibrium; kontribusi individu \\
~\cite{di_coupling_2023} & Simulator evolutionary game & Tingkat kerja sama; fraksi kooperator \\
~\cite{perera_learning_2025} & Repeated matrix games dengan populasi sintetis & Payoff rata-rata; robustness; generalisasi \\
~\cite{li_exploiting_2025} & Repeated zero-sum games & Galat prediksi; payoff kumulatif \\
~\cite{freire_modeling_2023} & Repeated matrix games & Efektivitas; stabilitas; akurasi prediksi \\
~\cite{hu_modeling_2023} & Simulasi repeated matrix game & Payoff rata-rata; kecepatan konvergensi \\
~\cite{elhamer_effects_2020} & Simulasi continuous-space & Tingkat kerja sama; stabilitas klaster \\
~\cite{wang_achieving_2019} & Lingkungan SPD 2D khusus & Reward rata-rata; kesejahteraan sosial \\
~\cite{jin_achieving_2025} & Benchmark MARL & Return ternormalisasi; konvergensi \\
~\cite{de_weerd_higher-order_2022} & Simulasi Colored Trails & Skor agen; kesejahteraan sosial \\
\midrule
\end{tabularx}
\end{table}

Sebagaimana terlihat pada Tabel~\ref{tab:evaluation_metrics}, mayoritas studi mengevaluasi efektivitas pemodelan lawan melalui metrik kinerja agregat jangka panjang, seperti payoff rata-rata, tingkat kerja sama, stabilitas, serta konvergensi menuju equilibrium. Evaluasi semacam ini menilai keberhasilan berdasarkan hasil akhir interaksi, tanpa secara eksplisit mengkaji bagaimana kualitas prediksi perilaku lawan pada berbagai horizon mempengaruhi proses pengambilan keputusan strategis.

Selain itu, sebagian besar pendekatan yang melibatkan komponen prediktif melaporkan kinerja berdasarkan galat satu langkah (\textit{one-step prediction error}). Meskipun metrik ini relevan untuk mengukur konsistensi lokal model terhadap data historis, ia tidak secara langsung mencerminkan kebutuhan perencanaan dalam \textit{repeated games}. Dalam konteks permainan berulang seperti IPD, keputusan pada suatu langkah mempengaruhi respons lawan pada langkah-langkah berikutnya, sehingga pengambilan keputusan bersifat intrinsik multi-langkah.

\begin{table}
\centering
\caption{Ringkasan Penelitian Terkait}\label{ringkasan_terkait}
\begin{tabularx}{\textwidth}{{p{1.5cm} p{3.5cm} p{8.0cm}}}
\toprule
\textbf{Ref} & \textbf{Nama Model} & \textbf{Temuan Utama} \\
\midrule

~\cite{qiao_o2m_2024} & Online Opponent Modeling (O2M) 
& O2M mampu beradaptasi lebih baik dan memperoleh rata-rata reward lebih tinggi dibandingkan baseline. \\

~\cite{zhu_evolutionary_2025} & Zero Determinant Strategy under Evolutionary Dynamic 
& Proporsi akhir kerja sama melebihi strategi Extort, menunjukkan keunggulan dinamika evolusioner. \\

~\cite{lv_inducing_2023} & Hierarchical Gifting DQN 
& Mendeteksi perubahan strategi lawan secara real-time dan dinamis menyesuaikan insentif kerja sama. \\

~\cite{gomez_grounded_2025} & Teamwork Game + MA-MAB 
& Mereproduksi pola mirip manusia (social loafing, compensation). \\

~\cite{di_coupling_2023} & Environment-Strategy Coupling Model 
& Kopling antara umpan balik lingkungan dan dinamika strategi secara signifikan meningkatkan stabilitas kerja sama kelompok. \\

~\cite{perera_learning_2025} & Ensemble-Training Cooperative Agent 
& Robust dan memiliki generalisasi lebih baik namun biaya komputasi lebih tinggi dan potensi estimasi reward intrinsik sub-optimal. \\

~\cite{li_exploiting_2025} & LSTM-Strategy 
& RNN mampu mengaproksimasi dinamika no-regret yang halus secara akurat serta memungkinkan eksploitasi yang menguntungkan pada tingkat non-stasioneritas rendah. \\

~\cite{freire_modeling_2023} & Internal Model 
& Konvergensi lebih cepat dan stabil dibanding TD-learning standar; performa menurun ketika perilaku lawan bergantung pada aksi agen sendiri (misalnya Tit-for-Tat). \\

~\cite{hu_modeling_2023} & Symmetric Learning Awareness (SLA) 
& Keseimbangan yang lebih stabil dan menghindari perilaku siklik pada metode gradien standar; pemodelan eksplisit meningkatkan konvergensi. \\

~\cite{elhamer_effects_2020} & Extended Social Particle Swarm (SPS) Model 
& Ukuran populasi dan laju pembaruan informasi membentuk pola kerja sama dinamis dan beragam. \\

~\cite{wang_achieving_2019} & Deep Multiagent Reinforcement Learning (untuk SPD) 
& Mencapai kerja sama mutual; adaptif terhadap perubahan strategi lawan. \\

~\cite{jin_achieving_2025} & Suggestion Sharing (SS) 
& Meningkatkan kerja sama dibanding value sharing dan policy sharing. \\

~\cite{de_weerd_higher-order_2022} & Higher-order Theory of Mind (ToM $k$, $k \geq 2$) 
& ToM tingkat tinggi berpengaruh signifikan; \textit{diminishing returns} setelah ToM-3. \\

\bottomrule
\end{tabularx}
\end{table}

Ketergantungan pada objektif prediksi satu langkah berpotensi menimbulkan ketidaksesuaian antara objektif statistik dan objektif strategis. Model yang dioptimalkan untuk meminimalkan galat satu langkah belum tentu menghasilkan estimasi lintasan perilaku lawan yang stabil ketika digunakan untuk perencanaan lintas-horizon. Dalam skema \textit{recursive forecasting}, kesalahan prediksi awal dapat terpropagasi dan mempengaruhi estimasi payoff kumulatif secara non-linear. Namun demikian, literatur yang ada jarang menguji secara eksplisit bagaimana kualitas prediksi multi-langkah berhubungan dengan performa strategis agen.

Lebih lanjut, belum banyak studi yang secara sistematis membandingkan implikasi penggunaan prediksi satu langkah, \textit{recursive multi-step forecasting}, dan \textit{direct multi-step forecasting} terhadap payoff strategis. Analisis mengenai keberadaan horizon prediksi optimal ($h^*$) yang memberikan trade-off terbaik antara stabilitas prediksi dan manfaat strategis juga masih terbatas.

Berdasarkan sintesis tersebut, terdapat celah penelitian yang signifikan, yaitu belum adanya analisis sistematis mengenai ketidaksesuaian antara objektif prediksi satu langkah dan kebutuhan perencanaan strategis multi-langkah dalam \textit{repeated games}. Secara khusus, hubungan antara akurasi prediksi lintas-horizon, propagasi kesalahan dalam skema rekursif, dan dampaknya terhadap payoff strategis belum dievaluasi secara terintegrasi dan terkontrol.

Penelitian ini berangkat dari celah tersebut dengan memfokuskan analisis pada evaluasi kinerja prediksi multi-langkah serta implikasinya terhadap pengambilan keputusan strategis dalam IPD, tanpa memosisikan arsitektur model sebagai kontribusi utama, melainkan sebagai instrumen untuk mengisolasi efek horizon peramalan.